% Section 3.2, from lectures/L03/appendix/b_interfaces.md.

\section{Interfaces in C++}
\label{c3:sec:interfaces}\label{c3:app:b}
An interface in C++ is an abstract base class that (normally) only contains pure virtual methods,
i.e.\ methods that have no implementation. The purpose of an interface is to define a common set
of methods that different classes can implement in their own way depending on their needs.

\subsection{Advantages of interfaces}
\label{c3:sec:advantages}
Using interfaces has several advantages:
\begin{itemize}
  \item \textbf{Abstraction:} You can write code that works against an abstraction rather than a
    concrete class. This makes it easier to replace implementations without changing the rest of
    the code.
  \item \textbf{Replaceability:} Different implementations can be used depending on the need, for
    example different hardware drivers or test classes.
  \item \textbf{Reusability:} Code that uses interfaces becomes more general and can be reused in
    different contexts.
  \item \textbf{Testability:} By using interfaces, you can easily create mock or fake classes for
    unit testing.
\end{itemize}

\subsection{Structure of an interface}
\label{c3:sec:structure}
Below the structure of an interface for timer circuits is shown. This interface is named
\code{driver::timer::Interface}. To keep the code compact, no comments are used in this example:

\begin{cppcode}
namespace driver::timer
{
class Interface
{
public:
    virtual ~Interface() noexcept = default;
    virtual void start() noexcept = 0;
    virtual void stop() noexcept = 0;
    [[nodiscard]] virtual bool isRunning() const noexcept = 0;
    virtual void reset() noexcept = 0;
};
} // namespace driver::timer
\end{cppcode}

\noindent Through the interface above, a timer can be started, stopped, and reset via the methods
\code{start()}, \code{stop()}, and \code{reset()}. It is also possible to check whether the timer
is active or not via the method \code{isRunning()}.

A few things to note:
\begin{itemize}
  \item The destructor is marked as \code{virtual}:
  \begin{itemize}
    \item This ensures that the correct destructor is called when an instance of a subclass is
      deleted through a pointer or reference to the interface.
    \item This is crucial for avoiding memory leaks and for correctly cleaning up resources in both
      the derived class and the base class.
  \end{itemize}
  \item The destructor is defined as \code{= default}:
  \begin{itemize}
    \item This results in an empty, automatic implementation.
    \item It is good practice to declare the destructor as virtual in an interface. Otherwise,
      objects of subclasses that are deleted through a pointer or reference to the interface may
      exhibit undefined behavior.
    \item Since the interface has nothing to clean up, a default implementation is sufficient here.
  \end{itemize}
  \item All methods that are meant to be overridden in subclasses are marked with \code{virtual}
    and end with \code{= 0}:
  \begin{itemize}
    \item This makes them pure virtual methods without implementation in the interface.
  \end{itemize}
  \item All virtual methods are here marked \code{noexcept}:
  \begin{itemize}
    \item This forces all concrete implementations in the subclasses to not be able to throw
      exceptions.
    \item This is advantageous in an embedded system, but if you do not want to enforce this it is
      fine to omit \code{noexcept} here.
  \end{itemize}
  \item The method \code{isRunning()} is marked \code{[[nodiscard]]}:
  \begin{itemize}
    \item This encourages callers to use the returned value by generating a compiler warning if it
      is discarded.
    \item The attribute belongs to the declaration it is written on, not to the class. It is
      therefore \textbf{not} inherited by overriding methods, and a call made directly on a
      concrete subclass whose implementation omits it produces no warning.
    \item Therefore, subclasses \textbf{should} repeat the \code{[[nodiscard]]} attribute on their
      implementations. This is done throughout this book.
    \item Note also that compiler support for \code{[[nodiscard]]} on virtual methods varies. Some
      compilers, such as GCC, do not warn at all when the return value of a \emph{virtual} call is
      discarded, not even when the call is made through the interface. Repeating the attribute on
      the override is what makes the warning appear where it can.
  \end{itemize}
\end{itemize}

\subsection{Structure of a concrete subclass}
\label{c3:sec:subclass}
Below is an example of a concrete subclass \code{driver::timer::Atmega328p}, which represents a
concrete implementation for timer circuits for the ATmega328P microcontroller:

\begin{cppcode}
namespace driver::timer
{
class Atmega328p final : public Interface
{
public:
    // Example constructor.
    explicit Atmega328p(std::uint16_t duration_ms) noexcept;

    // Overridden methods.
    ~Atmega328p() noexcept override;
    void start() noexcept override;
    void stop() noexcept override;
    [[nodiscard]] bool isRunning() const noexcept override;
    void reset() noexcept override;

    // Additional ATmega328P-specific methods (if any).

    // Deleted constructors and assignment operators.
    Atmega328p()                             = delete;
    Atmega328p(const Atmega328p&)            = delete;
    Atmega328p(Atmega328p&&)                 = delete;
    Atmega328p& operator=(const Atmega328p&) = delete;
    Atmega328p& operator=(Atmega328p&&)      = delete;

private:
    // Timer implementation details.
};
} // namespace driver::timer
\end{cppcode}

\noindent Some things we saw earlier in the section on inheritance (\secref{c3:sec:inheritance})
also appear here:
\begin{itemize}
  \item The keyword \code{final} is used to ensure that the class \code{driver::timer::Atmega328p}
    cannot be inherited; this is the final implementation:
  \begin{itemize}
    \item As a rule of thumb, this is good to do if the class is not intended to be inherited.
    \item However, in some cases it may be desirable to inherit the implementation for testing, and
      in that case this keyword must be omitted.
  \end{itemize}
  \item The class \code{driver::timer::Atmega328p} inherits the class
    \code{driver::timer::Interface} by writing \code{: public Interface} after the class name.
    Writing \code{public} means that everything that is public (or protected) in the interface is
    also public (or protected) in this derived class.
\end{itemize}
Some new things to note:
\begin{itemize}
  \item The overridden methods are marked with \code{override} precisely to indicate that these are
    concrete implementations of overridden virtual methods.
  \item The destructor is also marked with \code{override} to clearly show that it replaces the
    virtual destructor in the interface. It is the \code{virtual} destructor in the interface that
    ensures the correct destructor is called during polymorphic destruction; \code{override} makes
    the compiler check that this is the case, since it is an error if the base class destructor is
    not virtual.
  \item The constructor has nothing to do with the interface and is therefore neither marked
    \code{virtual} nor \code{override}.
\end{itemize}

\subsection{Example of using interfaces}
\label{c3:sec:usinginterfaces}
Below an interface \code{driver::led::Interface} is shown. This interface should be usable for
various LEDs, for example from different microprocessors:

\begin{cppcode}
/**
 * @file LED driver interface.
 */
#pragma once

#include <cstdint>

namespace driver::led
{
/**
 * @brief LED driver interface.
 */
class Interface
{
public:
    /**
     * @brief Destructor.
     */
    virtual ~Interface() noexcept = default;

    /**
     * @brief Get the pin the LED is connected to.
     *
     * @return The pin the LED is connected to.
     */
    [[nodiscard]] virtual std::uint8_t pin() const noexcept = 0;

    /**
     * @brief Check whether the LED is enabled.
     *
     * @return True if the LED is enabled, false otherwise.
     */
    [[nodiscard]] virtual bool isEnabled() const noexcept = 0;

    /**
     * @brief Enable/disable the LED.
     *
     * @param[in] enable True if the LED is to be enabled, false otherwise.
     */
    virtual void setEnabled(bool enable) noexcept = 0;

    /**
     * @brief Toggle the LED.
     */
    virtual void toggle() noexcept = 0;
};
} // namespace driver::led
\end{cppcode}

\noindent Below the subclass \code{driver::led::Atmega328p} is shown, which inherits the interface
\code{driver::led::Interface} and is implemented to easily control LEDs connected to the ATmega328P
microcontroller:

\begin{cppcode}
namespace driver::led
{
/**
 * @brief LED driver for ATmega328P.
 *
 *        This class is non-copyable and non-movable.
 */
class Atmega328p final : public Interface
{
public:
    /**
     * @brief Constructor.
     *
     * @param[in] pin The pin the LED is connected to.
     */
    explicit Atmega328p(std::uint8_t pin) noexcept;

    /**
     * @brief Destructor.
     */
    ~Atmega328p() noexcept override;

    /**
     * @brief Get the pin the LED is connected to.
     *
     * @return The pin the LED is connected to.
     */
    [[nodiscard]] std::uint8_t pin() const noexcept override;

    /**
     * @brief Check whether the LED is enabled.
     *
     * @return True if the LED is enabled, false otherwise.
     */
    [[nodiscard]] bool isEnabled() const noexcept override;

    /**
     * @brief Enable/disable the LED.
     *
     * @param[in] enable True if the LED is to be enabled, false otherwise.
     */
    void setEnabled(bool enable) noexcept override;

    /**
     * @brief Toggle the LED.
     */
    void toggle() noexcept override;

    Atmega328p()                             = delete; // No default constructor.
    Atmega328p(const Atmega328p&)            = delete; // No copy constructor.
    Atmega328p(Atmega328p&&)                 = delete; // No move constructor.
    Atmega328p& operator=(const Atmega328p&) = delete; // No copy assignment.
    Atmega328p& operator=(Atmega328p&&)      = delete; // No move assignment.

private:
    /** The pin the LED is connected to. */
    std::uint8_t myPin;

    /** Indicate whether the LED is enabled. */
    bool myIsEnabled;
};
} // namespace driver::led
\end{cppcode}

\noindent By using pointers or references to \code{driver::led::Interface}, you can write code that
works with all subclasses. As an example, a function named \code{blinkLed()} is demonstrated below,
which is used to blink a given LED. Note that:
\begin{itemize}
  \item The input argument \code{led} consists of a reference to a \code{driver::led::Interface}.
  \item The LED can therefore be an instance of any arbitrary subclass.
  \item For example, this could be an instance of the previously demonstrated class
    \code{driver::led::Atmega328p}, or alternatively it could be an instance of a class for an LED
    on an ESP32-S3 processor or similar.
\end{itemize}

\begin{cppcode}
/**
 * @brief Blink the given LED.
 *
 * @param[in, out] led The LED to blink.
 * @param[in] blinkTimeMs The blink time in milliseconds.
 */
void blinkLed(driver::led::Interface& led, const std::uint16_t blinkTimeMs) noexcept
{
    // Toggle the LED, then delay the calling thread.
    led.toggle();
    std::this_thread::sleep_for(std::chrono::milliseconds(blinkTimeMs));

    // Toggle the LED again, then delay the calling thread.
    led.toggle();
    std::this_thread::sleep_for(std::chrono::milliseconds(blinkTimeMs));
}
\end{cppcode}

\noindent When the function above is called, we can use one of the subclasses directly, for example
our class \code{driver::led::Atmega328p}. Assume that we have implemented an LED connected to pin 9
on an Arduino Uno via an instance named \code{led1}:

\begin{cppcode}
driver::led::Atmega328p led1{9U};
\end{cppcode}

\noindent We can blink this LED with a blink time of, for example, \code{1000 ms} by calling the
function \code{blinkLed()}. Since \code{driver::led::Atmega328p} is a subclass of
\code{driver::led::Interface}, we can pass \code{led1} directly:

\begin{cppcode}
blinkLed(led1, 1000U);
\end{cppcode}

\noindent Assume that we have also created a subclass \code{driver::led::Esp32s3} to implement LEDs
for an ESP32-S3 processor, as shown below.

Note that:
\begin{itemize}
  \item The implementation is slightly different this time; most notably, the user has the
    possibility to set the LED's initial state directly via a call to the constructor.
  \item This class also has a private method named \code{init()}.
  \item This is demonstrated to show that subclasses can be tailored according to specific needs:
\end{itemize}

\begin{cppcode}
namespace driver::led
{
/**
 * @brief LED driver for ESP32-S3.
 *
 *        This class is non-copyable and non-movable.
 */
class Esp32s3 final : public Interface
{
public:
    /**
     * @brief Constructor.
     *
     * @param[in] pin The pin the LED is connected to.
     * @param[in] initialState Initial state of the LED (default = off).
     */
    explicit Esp32s3(std::uint8_t pin, bool initialState = false) noexcept;

    /**
     * @brief Destructor.
     */
    ~Esp32s3() noexcept override;

    /**
     * @brief Get the pin the LED is connected to.
     *
     * @return The pin the LED is connected to.
     */
    [[nodiscard]] std::uint8_t pin() const noexcept override;

    /**
     * @brief Check whether the LED is enabled.
     *
     * @return True if the LED is enabled, false otherwise.
     */
    [[nodiscard]] bool isEnabled() const noexcept override;

    /**
     * @brief Enable/disable the LED.
     *
     * @param[in] enable True if the LED is to be enabled, false otherwise.
     */
    void setEnabled(bool enable) noexcept override;

    /**
     * @brief Toggle the LED.
     */
    void toggle() noexcept override;

    Esp32s3()                          = delete; // No default constructor.
    Esp32s3(const Esp32s3&)            = delete; // No copy constructor.
    Esp32s3(Esp32s3&&)                 = delete; // No move constructor.
    Esp32s3& operator=(const Esp32s3&) = delete; // No copy assignment.
    Esp32s3& operator=(Esp32s3&&)      = delete; // No move assignment.

private:
    /**
     * @brief Initialize the LED.
     */
    void init(bool initialState) noexcept;

    /** The pin the LED is connected to. */
    std::uint8_t myPin;

    /** Indicate whether the LED is enabled. */
    bool myIsEnabled;
};
} // namespace driver::led
\end{cppcode}

\noindent Assume that we have implemented an LED connected to pin 20 on an ESP32-S3 processor via an
instance named \code{led2}. We turn on the LED directly at startup:

\begin{cppcode}
driver::led::Esp32s3 led2{20U, true};
\end{cppcode}

\noindent We can also blink this LED by calling the function \code{blinkLed()}, since the class
\code{driver::led::Esp32s3} is a subclass of \code{driver::led::Interface}. For example, to blink
this LED with a blink time of 500 milliseconds, the following call can be made:

\begin{cppcode}
blinkLed(led2, 500U);
\end{cppcode}

\noindent The full example is in \secref{c3:sec:example}.

See also the following C implementation to understand how interfaces work ``under the hood'' (from
the previous embedded C course):
\url{https://github.com/qrtech-academy/embedded-c/blob/main/lectures/L10/interface_demo/README.md}.
